% GENERATED FILE. Edit canonical lessons, then run npm run generate:books.
% canonical-source-hash: fnv1a64:41f9c83931a70a96
% canonical-lessons: KA-C200-gunavagu, KA-C200-turisu, KA-C200-sayu, KA-C200-mugisu, KA-C200-suru-madu

\chapter{To get well, To itch, To die, To finish, To start}
\label{ch:ka-to-get-well-to-itch-to-die-to-finish-to-start}
\hlchaptermodality{200}

\begin{chapteropening}
\textbf{By the end of this chapter:} \emph{I can say to get well, to itch, to die, to finish and to start.}

Complete the last lesson of chapter 200: I can say to get well, to itch, to die, to finish and to start.
\end{chapteropening}

\section[\texorpdfstring{\kn{ಗುಣವಾಗು} (guṇavāgu)}{ಗುಣವಾಗು (guṇavāgu)}]{\kn{ಗುಣವಾಗು} (guṇavāgu) — to get well, to heal}

\label{lesson:KA-C200-gunavagu}



\begin{quote}
Before the new one: say the Kannada for to iron, then the Kannada for to change.
\end{quote}

\subsection*{You'll want to know: \kn{ಗುಣವಾಗು}}
\textbf{\kn{ಗುಣವಾಗು}} — \emph{guṇavāgu} — \textquotedblleft{}to get well, to heal\textquotedblright{}.

\begin{grammarlens}[title={What you've built}]
One more word for this chapter.
\end{grammarlens}

\subsection*{Your turn}
\begin{itemize}
\raggedright
  \item \emph{Say it:} \emph{guṇavāgu}
  \item \emph{Say it:} \emph{guṇavāgu}, once more
  \item \emph{Say it:} say \emph{guṇavāgu}
  \item \emph{Recall:} say the Kannada for to iron, then the Kannada for to change, then say \emph{guṇavāgu} again
\end{itemize}

\begin{tcolorbox}[breakable,colback=teal!4,colframe=teal!35!black,title={Before you move on}]
What does \kn{ಗುಣವಾಗು} mean? (\textquotedblleft{}To get well, to heal\textquotedblright{}.) Say it once more.
\end{tcolorbox}
\section[\texorpdfstring{\kn{ತುರಿಸು} (turisu)}{ತುರಿಸು (turisu)}]{\kn{ತುರಿಸು} (turisu) — to itch}

\label{lesson:KA-C200-turisu}



\begin{quote}
Before the new one: say the Kannada for to change, then the Kannada for to get well.
\end{quote}

\subsection*{You'll want to know: \kn{ತುರಿಸು}}
\textbf{\kn{ತುರಿಸು}} — \emph{turisu} — \textquotedblleft{}to itch\textquotedblright{}.

\begin{grammarlens}[title={What you've built}]
One more word for this chapter.
\end{grammarlens}

\subsection*{Your turn}
\begin{itemize}
\raggedright
  \item \emph{Say it:} \emph{turisu}
  \item \emph{Say it:} \emph{turisu}, once more
  \item \emph{Say it:} say \emph{turisu}
  \item \emph{Recall:} say the Kannada for to change, then the Kannada for to get well, then say \emph{turisu} again
\end{itemize}

\begin{tcolorbox}[breakable,colback=teal!4,colframe=teal!35!black,title={Before you move on}]
What does \kn{ತುರಿಸು} mean? (\textquotedblleft{}To itch\textquotedblright{}.) Say it once more.
\end{tcolorbox}
\section[\texorpdfstring{\kn{ಸಾಯು} (sāyu)}{ಸಾಯು (sāyu)}]{\kn{ಸಾಯು} (sāyu) — to die}

\label{lesson:KA-C200-sayu}



\begin{quote}
Before the new one: say the Kannada for to get well, then the Kannada for to itch.
\end{quote}

\subsection*{You'll want to know: \kn{ಸಾಯು}}
\textbf{\kn{ಸಾಯು}} — \emph{sāyu} — \textquotedblleft{}to die\textquotedblright{}.

\begin{grammarlens}[title={What you've built}]
One more word for this chapter.
\end{grammarlens}

\subsection*{Your turn}
\begin{itemize}
\raggedright
  \item \emph{Say it:} \emph{sāyu}
  \item \emph{Say it:} \emph{sāyu}, once more
  \item \emph{Say it:} say \emph{sāyu}
  \item \emph{Recall:} say the Kannada for to get well, then the Kannada for to itch, then say \emph{sāyu} again
\end{itemize}

\begin{tcolorbox}[breakable,colback=teal!4,colframe=teal!35!black,title={Before you move on}]
What does \kn{ಸಾಯು} mean? (\textquotedblleft{}To die\textquotedblright{}.) Say it once more.
\end{tcolorbox}
\section[\texorpdfstring{\kn{ಮುಗಿಸು} (mugisu)}{ಮುಗಿಸು (mugisu)}]{\kn{ಮುಗಿಸು} (mugisu) — to finish}

\label{lesson:KA-C200-mugisu}



\begin{quote}
Before the new one: say the Kannada for to itch, then the Kannada for to die.
\end{quote}

\subsection*{You'll want to know: \kn{ಮುಗಿಸು}}
\textbf{\kn{ಮುಗಿಸು}} — \emph{mugisu} — \textquotedblleft{}to finish\textquotedblright{}.

\begin{grammarlens}[title={What you've built}]
One more word for this chapter.
\end{grammarlens}

\subsection*{Your turn}
\begin{itemize}
\raggedright
  \item \emph{Say it:} \emph{mugisu}
  \item \emph{Say it:} \emph{mugisu}, once more
  \item \emph{Say it:} say \emph{mugisu}
  \item \emph{Recall:} say the Kannada for to itch, then the Kannada for to die, then say \emph{mugisu} again
\end{itemize}

\begin{tcolorbox}[breakable,colback=teal!4,colframe=teal!35!black,title={Before you move on}]
What does \kn{ಮುಗಿಸು} mean? (\textquotedblleft{}To finish\textquotedblright{}.) Say it once more.
\end{tcolorbox}
\section[\texorpdfstring{\kn{ಶುರು} \kn{ಮಾಡು} (śuru māḍu)}{ಶುರು ಮಾಡು (śuru māḍu)}]{\kn{ಶುರು} \kn{ಮಾಡು} (śuru māḍu) — to start}

\label{lesson:KA-C200-suru-madu}



\begin{quote}
Before the new one: say the Kannada for to die, then the Kannada for to finish.
\end{quote}

\subsection*{You'll want to know: \kn{ಶುರು} \kn{ಮಾಡು}}
\textbf{\kn{ಶುರು} \kn{ಮಾಡು}} — \emph{śuru māḍu} — \textquotedblleft{}to start\textquotedblright{}.

\begin{grammarlens}[title={What you've built}]
One more word for this chapter.
\end{grammarlens}

\subsection*{Your turn}
\begin{itemize}
\raggedright
  \item \emph{Say it:} \emph{śuru māḍu}
  \item \emph{Say it:} \emph{śuru māḍu}, once more
  \item \emph{Say it:} say \emph{śuru māḍu}
  \item \emph{Recall:} say the Kannada for to die, then the Kannada for to finish, then say \emph{śuru māḍu} again
\end{itemize}

\begin{tcolorbox}[breakable,colback=teal!4,colframe=teal!35!black,title={Before you move on}]
What does \kn{ಶುರು} \kn{ಮಾಡು} mean? (\textquotedblleft{}To start\textquotedblright{}.) Say it once more.
\end{tcolorbox}
